% SYNTHETIC TEST UNIT (integration tools) -- deliberate errors are marked "DELIBERATE"
\chapter{深度非弹散射（合成测试）}\label{ch:dis}

\section{运动学}\label{sec:dis-kin}

\begin{definition}[DIS 运动学]\label{def:dis-kinematics}
合成测试。
\uses{def:intro-tags}
\end{definition}

\begin{theorem}[欧氏 OPE \tagA]\label{thm:dis-ope}
合成测试。（本定理只涉及部分子态。）
% errtype: E-S
\uses{def:dis-kinematics}
\end{theorem}
\begin{proofsketch}略。\end{proofsketch}

\paragraph{地位与缺口}全阶论证；文献\cite{Collins2011}。

\begin{theorem}[矩的因子化 \tagR]\label{thm:dis-moments}
（$R_{\mathrm{OPE}}$：以 OPE 为前提而严格。QCD 中含夸克的微扰 OPE 是全阶论证（\cref{thm:dis-ope}），严格证明只对标量理论与纯 Yang--Mills 的规范不变算符给出；对强子态还须非微扰 OPE（\cref{asm:intro-np-ope}）。）
合成测试。
非微扰地位：\cref{asm:intro-wightman}、\cref{asm:intro-np-ope}（MA0）。
\uses{def:dis-kinematics,thm:dis-ope,asm:intro-wightman,asm:intro-np-ope}
\end{theorem}

\begin{proposition}[矩的推论 \tagR]\label{prop:dis-short}
（$R_{\mathrm{OPE}}$：以 OPE 为前提。）合成测试。（本命题只涉及部分子态。）
\uses{thm:dis-moments}
\end{proposition}

\begin{lemma}[微扰引理 \tagR]\label{lem:dis-ptlemma}
（$R_{\mathrm{PT}}$：在重整化微扰论中对红外正规化的部分子态严格。）合成测试。（本引理只涉及部分子态。）
\uses{def:dis-kinematics}
\end{lemma}

% DELIBERATE C-P3/K4 forward dependency on Ch12; C-P5 [R] uses an [A] statement
\begin{lemma}[前向依赖的引理 \tagR]\label{lem:dis-forward}
（$R_{\mathrm{PT}}$：在重整化微扰论中对红外正规化的部分子态严格。）合成测试。（本引理只涉及部分子态。）
\uses{thm:dy-factorization,thm:dis-ope}
\end{lemma}

% DELIBERATE C-P7 R_def uses an R_PT statement
\begin{proposition}[定义推论 \tagR]\label{prop:dis-rdef}
（$R_{\mathrm{def}}$：由算符定义直接推出。）合成测试。（本命题只涉及部分子态。）
\uses{lem:dis-ptlemma}
\end{proposition}

% DELIBERATE C-P7/K6 qualifier sentence not verbatim (\rm instead of \mathrm)
\begin{proposition}[限定语不逐字 \tagR]\label{prop:dis-notverbatim}
（$R_{\rm PT}$：在重整化微扰论中对红外正规化的部分子态严格。）合成测试。（本命题只涉及部分子态。）
\uses{def:dis-kinematics}
\end{proposition}

\begin{proposition}[被破坏的命题 \tagV]\label{prop:dis-violated}
合成测试：显式反例。（本命题只涉及部分子态。）
\uses{def:dis-kinematics}
\end{proposition}

\paragraph{地位与缺口}[V]：显式单圈反例\cite{Collins2011}。

% DELIBERATE C-P6 a positive statement uses a [V] statement
\begin{theorem}[使用 V 的定理 \tagA]\label{thm:dis-usesv}
合成测试。（本定理只涉及部分子态。）
\uses{prop:dis-violated,thm:dis-ope}
\end{theorem}

\paragraph{地位与缺口}全阶论证，瓶颈为\cref{thm:dis-ope}。

\begin{example}[例子]\label{ex:dis-example}
合成测试例子。
\uses{thm:dis-moments}
\end{example}

% DELIBERATE C-P1 W tagged remark without \uses
\begin{remark}[带标签的注 \tagA]\label{rem:dis-tagged}
合成测试。
\end{remark}

\paragraph{地位与缺口}全阶论证。

% DELIBERATE C-P1/K13 E two \uses lines
\begin{lemma}[两个依赖行 \tagR]\label{lem:dis-twouses}
（$R_{\mathrm{PT}}$：在重整化微扰论中对红外正规化的部分子态严格。）合成测试。（本引理只涉及部分子态。）
\uses{def:dis-kinematics}
\uses{lem:dis-ptlemma}
\end{lemma}

% DELIBERATE C-P1/K13 E no \uses line
\begin{theorem}[缺依赖行 \tagR]\label{thm:dis-nouses}
（$R_{\mathrm{PT}}$：在重整化微扰论中对红外正规化的部分子态严格。）合成测试。（本定理只涉及部分子态。）
\end{theorem}

% DELIBERATE C-P0 W optional title with unbraced nested brackets
\begin{lemma}[区间 [0,1] 上的界]\label{lem:dis-brackets}
合成测试（纯数学）。
% uses: none
\end{lemma}

% DELIBERATE C-P0 W \label not within the first three lines
\begin{definition}[标签位置]
第一行。

第三行。
\label{def:dis-latelabel}
% uses: none
\end{definition}

% DELIBERATE C-P0 W unlabelled environment
\begin{remark}
没有标签的注（合成测试）。
\end{remark}

% DELIBERATE C-P2/K5 W weakened (registry [A], written [P])
\begin{proposition}[被减弱的命题 \tagP]\label{prop:dis-weakened}
合成测试。（本命题只涉及部分子态。）
\uses{def:dis-kinematics}
\end{proposition}

\paragraph{地位与缺口}部分：只在单圈证明\cite{Collins2011}。

% DELIBERATE C-P2/K5 E strengthened (registry [P], written [A])
\begin{proposition}[被加强的命题 \tagA]\label{prop:dis-strengthened}
合成测试。（本命题只涉及部分子态。）
\uses{def:dis-kinematics}
\end{proposition}

\paragraph{地位与缺口}全阶论证。

% DELIBERATE C-P1 E \uses names a section label
\begin{example}[节标签依赖]\label{ex:dis-seclabel}
合成测试。
\uses{sec:dis-kin}
\end{example}

% DELIBERATE C-P3 W hidden forward dependency (through an appendix lemma that uses Ch12)
\begin{lemma}[隐藏的前向依赖 \tagR]\label{lem:dis-hidden}
（$R_{\mathrm{PT}}$：在重整化微扰论中对红外正规化的部分子态严格。）合成测试。（本引理只涉及部分子态。）
\uses{lem:appN-hidden}
\end{lemma}

% DELIBERATE K12 label list with a space (prints ??; figsync --normalize-refs repairs the copy);
% K2 I canonical label of the absent unit Ch05; K15 unknown citation key
见\cref{thm:dis-ope, thm:dis-moments}与\cref{thm:sudakov-factorization}；文献\cite{Synth:2026abc,NoSuchKey2099}。

%%% BEGIN fig:synth-same
\begin{figure}[htbp]
\centering
\begin{tikzpicture}[fd]
\draw[quark] (0,0) -- (2,0);
\end{tikzpicture}
\caption{合成测试图甲（与单元中的副本一致）。}\label{fig:synth-same}
\end{figure}
%%% END fig:synth-same

% DELIBERATE figsync DIFFERS: old caption (the source has been corrected since)
%%% BEGIN fig:synth-changed
\begin{figure}[htbp]
\centering
\begin{tikzpicture}[fd]
\draw[gluon] (0,0) -- (2,0);
\end{tikzpicture}
\caption{合成测试图乙（旧图题）。}\label{fig:synth-changed}
\end{figure}
%%% END fig:synth-changed

%%% BEGIN fig:synth-existing
\begin{figure}[htbp]
\centering
\begin{tikzpicture}[fd]
\draw[photon] (0,0) -- (2,0);
\end{tikzpicture}
\caption{合成测试图丙（已有图）。}\label{fig:synth-existing}
\end{figure}
%%% END fig:synth-existing
